% Part: computability
% Chapter: computability-theory
% Section: ce-closed-cup-cap

\documentclass[../../../include/open-logic-section]{subfiles}

\begin{document}

\olfileid{cmp}{thy}{clo}

\olsection[గ.లె. సమితుల సంయోగం, ఛేదనం]{గణనీయంగా లెక్కించదగిన
  సమితులు సంయోగం, ఛేదనం కింద సంవృతాలు}

గణనీయంగా లెక్కించదగిన సమితుల వర్గానికి గల కొన్ని సంవృతతా ధర్మాలను
కింది సిద్ధాంతం ఇస్తుంది.

\begin{thm}
$A$, $B$లు గణనీయంగా లెక్కించదగినవని అనుకుందాం. అప్పుడు $A \cap B$,
$A \cup B$లు కూడా గణనీయంగా లెక్కించదగినవే.
\end{thm}

\begin{proof}
గణనీయంగా లెక్కించదగిన సమితులకు గల వివిధ లక్షణీకరణలను వాడటానికి
\olref[eqc]{thm:ce-equiv} అనుమతిస్తుంది. వివరణ కోసం కొన్ని భిన్న
నిరూపణలను ఇస్తాం.

మొదటి నిరూపణలో, $A$ను గణనీయ ప్రమేయం~$f$, $B$ను గణనీయ
ప్రమేయం~$g$ లెక్కిస్తాయని అనుకుందాం. ఇలా ఉంచండి:
\begin{align*}
h(x) & = \umin{y}{(f(y) = x \lor g(y) = x)} \text{ మరియు}\\
j(x) & = \umin{y}{(f((y)_0) = x \land g((y)_1) = x)}.
\end{align*}
అప్పుడు $A \cup B$ అనేది $h$ యొక్క నిర్వచన క్షేత్రం; $A \cap B$
అనేది~$j$ యొక్క నిర్వచన క్షేత్రం.

\begin{explain}
గణనా పరంగా ఇక్కడ జరుగుతున్నది ఇదే: $A$, $B$లను లెక్కించే ప్రక్రియలతో
పాటు మూలకం~$x$ ఇచ్చినప్పుడు, రెండు లెక్కింపుల్లో దేనిలోనైనా~$x$ కోసం
వెతికి అది $A \cup B$లో ఉందో అర్ధ-నిర్ణయించవచ్చు. అలాగే మూలకం~$x$
అనేది $A \cap B$లో ఉందో అర్ధ-నిర్ణయించడానికి, రెండు లెక్కింపుల్లో
ఒకేసారి~$x$ కోసం వెతకవచ్చు.
\end{explain}

రెండవ నిరూపణలోనూ $A$ను $f$, $B$ను $g$ లెక్కిస్తాయని అనుకుందాం.
ఇలా ఉంచండి:
\[
k(x) = \begin{cases}
f(x/2) & \text{$x$ సరిసంఖ్య అయితే} \\
g((x-1)/2) & \text{$x$ బేసి సంఖ్య అయితే.}
\end{cases}
\]
అప్పుడు $k$, $A \cup B$ను లెక్కిస్తుంది; $f$,~$g$ ఇచ్చే లెక్కింపులను
$k$ వంతులవారీగా నడుపుతుందన్నదే భావన. $A \cap B$ను లెక్కించడం
కొంచెం క్లిష్టం. $A \cap B$ శూన్యమైతే, నిర్వచనం ప్రకారమే అది
గణనీయంగా లెక్కించదగినది. లేకపోతే $A \cap B$లో ఏదైనా మూలకం $c$ను
తీసుకుని, $l$ను ఇలా నిర్వచించండి:
\[
l(x) = \begin{cases}
f((x)_0) & \text{$f((x)_0) = g((x)_1)$ అయితే} \\
c & \text{ఇతర సందర్భాల్లో.}
\end{cases}
\]
గణనా పరంగా, $l$ అనేది $f$, $g$ల లెక్కింపుల్లోని మూలకాల జతలన్నిటినీ
పరిశీలించి, తనకు దొరికిన ప్రతి సరిపోలికను అవుట్‌పుట్ చేస్తుంది;
సరిపోలిక లేనప్పుడు కొత్త విలువను ఇవ్వకుండా $c$నే అవుట్‌పుట్ చేస్తుంది.
\sourcecorrection{OLTECOMTHY-007}{ఈ నిరూపణలోని ఆంగ్ల మూలంలో ఒకచోట ``by looking''కు బదులు ``for looking'' అని, మరొకచోట ``trickier''కు బదులు ``tricker'' అని ఉంది. వాదన కోరే సాధారణ అర్థాన్ని తెలుగులో సరిగా ఇచ్చాం.}

చివరి నిరూపణలో, $A$ అనేది పాక్షిక ప్రమేయం $m(x)$ యొక్క
\emph{నిర్వచన క్షేత్రం}, $B$ అనేది పాక్షిక ప్రమేయం $n(x)$ యొక్క
నిర్వచన క్షేత్రం అనుకుందాం. అప్పుడు $A \cap B$ అనేది పాక్షిక
ప్రమేయం $m(x)+n(x)$ యొక్క నిర్వచన క్షేత్రం.

\begin{explain}
గణనా పరంగా, $m$ ఆగే విలువల సమితి $A$, $n$ ఆగే విలువల సమితి $B$
అయితే, రెండు ప్రక్రియలూ ఆగే విలువల సమితియే $A \cap B$.
\end{explain}

$A \cup B$ను ఆగే విలువల సమితిగా వ్యక్తపరచడం మరింత క్లిష్టం;
ఎందుకంటే $m$, $n$లను సమాంతరంగా అనుకరించాలి. $d$ అనేది $m$కు
సూచిక, $e$ అనేది $n$కు సూచిక అనుకుందాం; అంటే $m=\cfind{d}$,
$n=\cfind{e}$. అప్పుడు $A \cup B$ కింది ప్రమేయపు నిర్వచన క్షేత్రం:
\[
p(x) = \umin{y}{(T(d,x,y) \lor T(e,x,y))}.
\]
\begin{explain}
గణనా పరంగా, ఇన్‌పుట్~$x$పై $p$, $m$కు లేదా $n$కు చెందిన ఆగే గణన
కోసం వెతికి, వాటిలో ఏదో ఒకటి దొరికితే ఆగుతుంది.
\end{explain}
\end{proof}

\end{document}
